% !TEX root = ../notes.tex

\documentclass[../notes.tex]{subfiles}

\begin{document}

\section{March 16}
Today, we continue with the hydrogen atom.

\subsection{The Eigenfunctions}
Let's start by finding eigenfunctions for the spherical Laplacian.
\begin{lemma}
	We compute the eigenvalues $\lambda$ and functions $\xi\in L^2\left(S^2\right)$ such that
	\[\Delta_{\mathrm{sph}}\xi=\lambda\xi.\]
\end{lemma}
\begin{proof}
	We are hunting for a function $\xi\in L^2\left(S^2\right)$ and $\lambda\in\CC$ for which $\Delta_{\mathrm{sph}}\xi=\lambda\xi$. Note that $\Delta_{\mathrm{sph}}$ is $\op{SO}(3)$-invariant, so the set of solutions will be a representation of $\op{SO}(3)$. By \Cref{thm:peter-weyl}, we know that
	\[L^2(\op{SO}(3))=\widehat{\bigoplus_m}L_{2m}\otimes L_{2m}^*.\]
	Here, $L_{2m}$ is the $2m$th symmetric power of the spin representation, and these are all the irreducible representations of $\op{SO}(3)$, which we see because $\op{Spin}(3)=\op{SL}(2)$. Note that $S^2=\op{SO}(3)/\op{SO}(2)$ as a homogeneous space, so it follows that
	\[L^2\left(S^2\right)=\widehat{\bigoplus_m}L_{2m}\otimes (L_{2m}^*)^{\op{SO}(2)},\]
	which one can calculate is
	\[L^2\left(S^2\right)=\widehat{\bigoplus_m}L_{2m}.\]
	For example, the algebraic functions $L^2_{\mathrm{alg}}\left(S^2\right)$ (namely, the polynomial functions on the sphere) is $\bigoplus_mL_{2m}$. Now, $\Delta_{\mathrm{sph}}$ can be checked to send trigonometric functions to trigonometric functions, so it descends to an $\operatorname{SO}(3)$-invariant function on $L^2_{\mathrm{alg}}\left(S^2\right)$, which must act by some eigenvalue $-\lambda_m$ on each $L_{2m}$.

	For example, let's hunt for the weight-$0$ eigenfunctions in $L_\ell$. Then they are invariant under $\op{SO}(2)$, so they should depend only on $\varphi$, so we can write them as some polynomial $Y^0_\ell=P_\ell(\cos\varphi)$ of degree $\ell$. By orthogonality of the matrix coefficients, one sees that
	\[\int_0^\pi P_\ell(\cos\varphi)P_k(\cos\varphi)\sin\varphi\,d\varphi=0\]
	for $\ell\ne k$. (Here, the $\sin\varphi$ comes from our choice of measure.) A change of coordinates makes this
	\[\int_{-1}^1P_\ell(z)P_k(z)\,dz=0\]
	for $\ell\ne k$, so we are looking at the Legendre polynomials! Let's compute the eigenvalues: because $P_\ell$ does not depend on $\theta$, we see that
	\[\Delta_{\mathrm{sph}}P=\del_z\left(1-z^2\right)\del_zP_\ell.\]
	If the eigenvalue is $-\lambda_\ell$, then a calculation on the leading term reveals that $\lambda_\ell=\ell(\ell+1)$.

	In general, for each $L_{2\ell}$, there is an eigenfunction $Y_\ell^m(\varphi,\theta)$ of weight $m\in[-\ell,\ell]$. Accordingly, $e^{-im\theta}Y_\ell^m$ is $\op{SO}(2)$-invariant,\todo{} so it takes the form $e^{im\theta}P_\ell^m(z)$. For even $m$, $P_\ell^m$ turns out to be a polynomial in $z$; otherwise, it is $\sqrt{1-z^2}$ times a polynomial in $z$. Our eigenvalue calculation from earlier reveals that
	\[\Delta_{\mathrm{sph}}Y_\ell^m=-\ell(\ell+1)Y_\ell^m,\]
	so one finds that $P_\ell^m$ satisfies the Legendre differential equation
	\[\del_z\left(1-z^2\right)\del_zP-\frac{m^2}{1-z^2}P+\ell(\ell+1)P=0.\]
	One can solve this differential equation.
\end{proof}
On the homework, we will solve this differential equation, and we will show that they take the following form.
\begin{definition}[generalized Legendre polynomials]
	For a nonnegative integer $\ell$ and $m\in[-\ell,\ell]$, the \textit{generalized Legendre polynomials} are
	\[P_\ell^m(z)=\left(1-z^2\right)^{m/2}\del_z^{\ell+m}\left(1-z^2\right)^\ell.\]
\end{definition}
\begin{remark}
	Here is a high-level explanation for these eigenvalues. There are angular momentum differential operators $L_x$, $L_y$, and $L_z$, which ``quantize'' the components of angular momentum (computed as $p\times r$). It turns out that $\{-iL_x,-iL_y,-iL_z\}$ spans $\mf{so}_3(\RR)$, and it turns out that the Casimir element
	\[L_x^2+L_y^2+L_z^2\]
	is $-\Delta_{\mathrm{sph}}$. However, one can calculate that the action of the Casimir on $L_{2\ell}$ is $\ell(\ell+1)$, which explains why we received those eigenvalues.
\end{remark}
We now turn to the part depending on $r$.
\begin{lemma}
	We compute the functions $f$ on $\RR$ and energies $E<0$ for which
	\[f''(r)+\frac2rf'(r)+\left(\frac2r-\frac\lambda{r^2}+2E\right)f(r)=0.\]
\end{lemma}
\begin{proof}
	It will be convenient to rescale $f$ to a function $h$ satisfying $f(r)=r^\ell e^{-r/n}h(2r/n)$, where $n$ is some positive real number. (It will later turn out that $n$ is a positive integer.) Applying this change of variables yields
	\[\rho h''(\rho)+(2\ell+2-\rho)h'(\rho)+\left(n-\ell-1+\frac{1+2En^2}4\rho\right)h(\rho)=0.\]
	Thus, it is convenient to set $n$ to be $\sqrt{-1/(2E)}$ so that the equation reads
	\[\rho h''(\rho)+(2\ell+2-\rho)h'(\rho)+(n-\ell-1)h(\rho)=0.\]
	This is called the generalized Laguerre equation. Notably, the condition $\norm\psi^2<\infty$ translates to
	\[\int_{\RR^+}\rho^{2\ell+2}e^{-\rho}\left|h(\rho)\right|^2\,d\rho<\infty.\]
	For example, we should check that we do not blow up at zero: if $h$ around zero looks like $\rho^s(1+o(\rho))$, then we find a characteristic equation $s(s+2\ell+1)=0$ (from the differential equation). Thus, $s=0$ or $s=-2\ell-1$. Let's explain how to throw out the latter case.
	\begin{itemize}
		\item With $\ell>0$, one finds that $\rho^{2\ell+2}\left|h(\rho)\right|^2\sim\rho^{-2\ell}$ if $s=-2\ell-1$ is not integrable at zero, so we may throw out this case.
		\item If $\ell=0$, then we can try to take $s=-2\ell-1=-1$. Then $\rho^2\left|h(\rho)\right|^2$ has finite limit at $\rho=0$, which appears to be integrable. In this case $\ell=0$, one finds that $\xi$ is constant, so $\psi$ depends only on the radius, and we see $\psi(r)\sim1/r$ near $r=0$. But then $H\psi-E\psi$ does not vanish as a distribution because it is roughly $\delta_0$. Thus, ew can throw out $s=0$ in this case as well.
	\end{itemize}
	Thus, for any $\ell$, we are required to work with $s=0$, producing a unique solution $h_n$ satisfying the initial condition $h_n(0)=1$. We may calculate $h_n$ using a power series. By substituting a power series, we find that
	\[h_n(\rho)=\sum_{k=0}^\infty\frac{(1+\ell-n)\cdots(k+\ell-n)}{(2\ell+2)\cdots(2\ell+1+k))}\cdot\frac{\rho^k}{k!}.\]
	The ratio test verifies that this converges for all $\rho$. If $n$ is not a positive integer, then computing the (large $k$) power series coefficients finds that $\left|h(\rho)\right|\sim e^\rho$, which is not possible by the $\norm\psi^2<\infty$ condition. On the other hand, $n$ being a positive integer makes $h(\rho)$ into a polynomial, which is successfully integrable.
\end{proof}
The moral is that our lemma has found that the solutions $h_n$ are the so-called generalized Laguerre polynomials.
\begin{definition}[generalized Laguerre polynomial]
	Choose a positive integer $N$ and real number $\alpha$. Then the \textit{generalized Laguerre polynomials} is the polynomials
	\[L^\alpha_N(\rho)\coloneqq\sum_{k=0}^N(-1)^N\frac{N\cdots(N-k+1)}{(\alpha+1)\cdots(\alpha+k)}\cdot\frac{\rho^k}{k!}.\]
\end{definition}
Combined with our discussion last class, we have proven the following theorem.
\begin{theorem} \label{thm:hydrogen-atom}
	The bound states of the hydrogen atom are, up to scaling,
	\[\psi_{n\ell m}=r^\ell e^{-r/n}L_{n-\ell-1}^{2\ell+1}\left(\frac{2r}n\right)Y_\ell^m(\theta,\varphi),\]
	where $Y_\ell^m(\theta,\varphi)=e^{im\theta}P_\ell^m(\varphi)$, and $n$ is a positive integer, and $\ell$ is a nonnegative integer at most $n-1$, and $m\in[-\ell,\ell]$. The energy level $E_n$ of this bound state is $E_n=-\frac1{2n^2}$.
\end{theorem}
Let's name the parameters for our quantum states.
\begin{defihelper}[quantum number] \nirindex{quantum number} \nirindex{quantum number!azimuthal quantum number} \nirindex{quantum number!magnetic quantum number}
	The \textit{quantum number} of the bound state $\psi_{n\ell m}$ is the number $n$. We also call $\ell$ the \textit{azimuthal quantum number}, and $m$ is the \textit{magnetic quantum number}.
\end{defihelper}
\begin{corollary}
	The space $W_n$ of bound states with principal quantum number $n$ is isomorphic to the representation
	\[L_0\oplus L_2\oplus\cdots\oplus L_{2n}\]
	of $\op{SO}(3)$. In particular, it has dimension $n^2$.
\end{corollary}
\begin{proof}
	This follows from the proof of \Cref{thm:hydrogen-atom}.
\end{proof}
\begin{remark}[spin]
	From the periodic table, one expects to have $2n^2$ states for a given $n$: for example, there are two places the electron can go to in the $s$ orbit (corresponding to the two electrons in the helium atom). It turns out that electrons have an extra degree of freedom, which is called the spin, taking values in $\{+1/2,-1/2\}$. Then the space of states lives in $L^2\left(\RR^3\right)\otimes\CC^2$. It turns out that $\CC^2$ experimentally looks like the two-dimensional representation of $\mf{so}(3)$, so the correct group acting should be $\op{Spin}(3)$ instead of $\op{SO}(3)$. This is not a contradiction in quantum mechanics (though it looks like a contradiction in classical physics) because physical states are only recovered from quantum ones up to a phase factor $e^{i\varepsilon}$, and the extra element in $\op{Spin}(3)$ only adjusts the quantum state up to a sign! Anyway, in the end, we get the correct $2n^2$ states.
\end{remark}

\subsection{Coulomb Waves}
The differential operator $\Delta_{\mathrm{sph}}$ fails to be compact, so we should not expect its spectrum to go to zero. Thus, it is surprising that we only found energy levels going to zero. The solution to this problem is that the bound states $\psi_{n\ell m}$ turn out to not be a basis of $L^2\left(\RR^3\right)$. Indeed, let $\varphi$ be some compactly supported function on $\RR^3$ at the origin. Then define $\varphi_s(x)\coloneqq\varphi(x+sa)$ to be some scaling.

We claim that $(H\varphi_s,\varphi_s)$ is positive for large $s$: indeed, by integration by parts, one can remove the Laplacian, and we find that
\[(H\varphi_s,H\varphi_s)=\frac12\int\left|\nabla\varphi\right|^2\,dV-\int\frac1{\norm x}\left|\varphi(x+sa)\right|^2\,dV.\]
By shifting coordinates, this is
\[(H\varphi_s,H\varphi_s)=\frac12\int\left|\nabla\varphi\right|^2\,dV-\int\frac{\left|\varphi(x+sa)\right|^2}{{\norm {x-sa}}}\,dV.\]
Now, as $s\to\infty$, the right integral goes to zero, and the left integral is constant and positive, so $(H\varphi_s,\varphi_s)$ is positive! Thus, $\varphi_s$ cannot be in the span of the $\psi_{n\ell m}$s for sign reasons.

Thus, it is possible to get positive energy levels.
\begin{definition}[Coulomb waves]
	The eigenfunctions with positive energy levels in $L^2\left(\RR^3\right)$ are called the \textit{Coulomb waves}.
\end{definition}
It turns out that the Coulomb waves turn out to have a continuous spectrum, but the core ideas used to understand them are rather similar to the bound states.

% \subsection{Spin}


\end{document}